\documentclass{article}
\usepackage[utf8]{inputenc}
\usepackage[T1]{fontenc}
\usepackage[margin=2cm]{geometry}
\usepackage{hyperref}
\sloppy
\begin{document}
\section*{Historical source}
Encyclopaedia Britannica scan, source row 2277. Parse 3.5 extracted the text. Super 3.5 VL corrected it from an enlarged crop of each block. It can still contain recognition errors.
Page headers, footers, and page numbers are removed.

\subsection*{ANATOMY (E0212)}
ANATOMY.

Of the muscles. 12 12, The first cervical nerves; 13 13, The second cervical nerves; both spent upon the muscles that lie on the neck, and teguments of the neck and head. 14 14, The third cervical nerves, which, after sending off (15 15, \&c.) the phrenic nerves to the diaphragm, supply the muscles and teguments that lie on the side of the neck and top of the shoulder. 16 16, The brachial plexus, formed by the fourth, fifth, sixth, seventh cervicals, and first dorsal nerves,—which supply the muscles, and teguments of the superior extremity. 17 17, The twelve dorsal, or proper intercostal nerves, which are spent upon the intercostal muscles and some of the large muscles which lie upon the thorax. 18 18, The five lumbar pairs of nerves, which supply the lumbar and abdominal muscles, and some of the teguments and muscles of the inferior extremity. 19 19, The sacro-sciatic, or posterior crural nerve, formed by the two inferior lumbar, and three superior of the os sacrum. This large nerve supplies the greatest part of the muscles and teguments of the inferior extremity. 20, The lumbosac plexus, formed by the eighth pair. 21 21, Branches of the solar or caulic plexus, formed by the eighth pair and inter-

coftals, which supply the stomach and chylopoietic
viscera. 22 22, Branches of the superior and inferior
mefenenteric plexufes, formed by the eighth pair and in-
tercoftals, which supply the chylopoietic viscera,
with part of the organs of urine and generation.
23 23, Nerves which accompany the spermatic cord.
24 24, The hypogastric plexus, which fuppiles the
organs of urine and generation within the pelviis.

Fig. 2, 3, 4, 5. Show different views of the Inferior part of the Brain, cut perpendicularly through the Middle—with the Origin and large portions of all the Nerves which pass out through the Bones of the Cranium,—and the three first Cervicals.

A, The anterior lobe. B, The lateral lobe of the cerebrum. C, One of the lobes of the cerebellum. D, Tuber annulare. E, Corpus pyramidale, in the middle of the medulla oblongata. F, The corpus olivare, in the side of the medulla oblongata. G, The medulla oblongata. H, The medulla spinalis.

Nerves.—1 2 3 4 5 6 7 8 and 9, Pairs of nerves.
10 10, Nervus accessorius, which comes from—11, 12,
and 13, the three first cervical nerves.

CHAP. VI. OF THE SENSES, AND THEIR ORGANS.

IN treating of the fenfes, we mean to confine ourselves to the external ones of touch, taft, smelling, hearing, and vision. The word sense, when applied to these five, feems to imply not only the sensation excited in the mind by certain impreffions made on the body, but likewise the organ defined to receive and transmit these impreffions to the sensorium. Each of these organs being of a peculiar structure, is susceptible only of particular impreffions, which will be pointed out as we proceed to describe each of them fe parately.

SECT. I. Of Touch.

141 In this feature of touch may be defined to be the faculty of diltinguling certain properties of bodies by the feel. In a general accepted, this definition might perhaps not properly be extended to every part of the body political of identifying (又), it is commonly combined to the nervous pupils, and consequently of exciting differences in the fenision communes. But the origin of touch, like all the other feels, it not equally distinct in every part of the body, or in every subjects, being in force much more equitable than it to others.

SECT. II. Of Taste.

The sense of taste is seated chiefly in the tongue ; the situation and figure of which are sufficiently known.

On the upper surface of this organ we may observe a great number of papillae, which, on account of their difference in size and shape, are commonly divided into three classes. The largest are situated towards the basis of the tongue. Their number commonly varies from seven to nine, and they seem to be mucous follicles. Those of the second class are somewhat smaller, and of a cylindrical shape. They are most numerous about the middle of the tongue. Those of the third class are very minute, and of a conical shape. They are

(r) In the course of this article, mention has often been made of the sensibility or insensibility of different parts of the body: it will therefore, perhaps, not be amiss to observe in this place, that many parts which were formerly supposed to possess the most exquisite sense, are now known to have but little or no feeling, at least in a sound state; for in an inflamed state, even the bones, the most insensible parts of any, become susceptible of the most painful sensations. This curious discovery is due to the late Baron Haller. His experiments prove, that the bones, cartilages, ligaments, tendons, epidermis, and membranes (as the pleura, pericardium, dura and pia mater, periosteum, \&c.), may in a healthy state be considered as insensible. As sensibility depends on the brain and nerves, of course different parts will possess a greater or less degree of feeling, in proportion as they are supplied with a greater or smaller number of nerves. Upon this principle it is, that the skin, muscles, intestines, urinary bladder, ureters, uterus, vagina, penis, tongue, and retina, are extremely sensible, while the lungs and glands have only an obscure degree of feeling.


\section*{Source record}
\begin{description}
\item[Scan ID] \texttt{britannica:309f765b1ecc83e0a7f7762e44b4a927a7555443:train:2277}
\item[Image file] \texttt{page-2277-cdde2fc730b9.jpg}
\item[Image SHA-256] \texttt{cdde2fc730b9b43e147b41385c61a207df3d0519a0ec368e9bbc3b3ecd4432de}
\item[Dataset] NationalLibraryOfScotland/encyclopaedia\_britannica\_illustrated at revision \texttt{309f765b1ecc83e0a7f7762e44b4a927a7555443} (CC0-1.0)
\item[Extraction] model \texttt{nvidia/nemotron-3.5-parse}, raw reply \texttt{extract/raw/page-2277-cdde2fc730b9.json}
\item[Raw reply SHA-256] \texttt{14347ce4d86a8b7c64b06fe9b4a47760c96bcdeeee7fd952886c483d168e7d17}
\item[Correction] model \texttt{nvidia/Nemotron-3.5-Super-VL}, block crops, raw replies \texttt{correct/raw/crop\_nothink/page-2277-cdde2fc730b9-b*.json}; parse\_fallback: 1, super\_vl: 13
\item[Running header] Chap. VI.
\end{description}
\end{document}
